% 本文件由 LaTeX 排版 Agent 根据有效 translation.json 自动生成。
% author=Novatian; work_id=0512; title=论犹太人的食物

\chapter{论犹太人的食物}

% structure: p0001 source=h1 target=omitted reason=page-title-already-rendered-by-parent

% structure: p0002 source=h2 target=omitted reason=duplicate-page-title

至圣的弟兄们，虽然我殷切盼望收到你们信件与著作的日子，并视之为最重要、最幸福的日子——因为如今还有什么能使我更加喜乐呢？——但我认为，当我以所欠你们的爱与相应的关切回复你们类似的书信时，那日子也不应视为等闲，而应视作特殊的日子。因为，至圣的弟兄们，没有什么能如此束缚我，也没有什么能以如此挂虑与焦虑的刺激使我激动不安，如同害怕你们会因我的缺席而认为你们遭受了任何损失；而我则努力补救，通过频繁的书信来表明我与你们同在。因此，虽然我所欠的责任、我所承担的使命、以及加于我身上的圣职，都要求我不得不写信，但你们仍通过你们不断的通信激励我写作，更加激发了这需要。虽然我虽已倾向于那些周期性的爱的表达，但你们通过表明你们恒常持守福音而更加催促我：由此导致的结果是，我的书信与其说是教导你们这些已经知情的人，不如说是激励你们这些已经预备好的人。因为你们不仅持守纯洁且洗净一切邪说污点的福音，还积极地教导这福音，你们不再寻求人为师傅，因为你们藉这些事表明自己是教师。因此，当你们奔跑时，我劝勉你们；当你们警醒时，我激励你们；当你们对抗\BibleQuote{邪恶的属灵事物}，\BibleRef{弗 6:12}时，我向你们说话；当你们\BibleQuote{向着基督召叫的奖赏竭力奔跑}，\BibleRef{斐 3:14}时，我催促你们——好叫你们践踏并拒绝异端者亵渎的诽谤以及犹太人无稽的传说，持守基督唯一的圣言和教导，从而配得祂名下之权威。但是，犹太人何等悖逆，对他们的法律之理解何等遥远，我相信我在前两封信中已充分表明，并且在那里已经确切证明他们不知道何为真割损、何为真安息日；而他们日益加增的盲目在本封书信中被驳斥，其中我简要论述了他们的食物，因为他们在这些食物上认为自己才是神圣的，而其他一切皆被玷污。\par

% structure: p0004 source=h2 target=section reason=relative-heading-rank; base=h2

\section{第二章　论述。——他首先断言法律是属神的；因此，人最初的食物只是果树，而后加上了肉食，以致随后而来的法律应被属神地理解。}

因此，首先，我们必须利用那段经文：\BibleQuote{法律是属神的}\BibleRef{罗 7:14}；他们若否认法律是属神的，便确然亵渎；若避免亵渎而宣认它是属神的，就当以属神的方式去读它。因为神圣的事物必须以神圣的方式接受，而确然应被持守为圣。但将世俗和人类的道理附会于神圣和属神的话语的人，被烙上严重的过犯；我们必须提防此事。此外，我们可以提防，若天主所命令的任何事被如此对待，好像被假定为削弱祂的权威，因称某些事物为不洁与不净，其设立反而使设立者蒙羞。因为在弃绝祂所造之物时，祂似乎要谴责那些祂曾嘉许为善的祂自己的工程；祂将被指为在两种情况下显得反复无常，正如异端者所愿：或曾祝福不洁净的事物，或后来弃绝为不善那些祂曾祝福为既洁净又善的受造物。若固执那犹太教义，此事的极大矛盾将永远存在；我们必须尽一切能力摆脱它，好使我们能除去他们不正规传授的一切，而恢复祂工程的合适安排，以及神圣法律的合宜而属神的应用。但是要从万物的起源开始，这也是我必须开始的：最初人类唯一的食物是果实和树木的出产。因为后来，人的罪将他的需要从果树转移到地上的出产，当时他身体的姿态本身便证明了他的良心状态。因为只要他们良心无愧，无罪便使他们昂首向天，从树上摘取食物；然而一旦犯罪，便使人俯向大地和地上去收取谷物。此后，又加上了肉的食用，天主眷顾为人提供合宜场合所需的各种肉类。因为需要一种更柔嫩的食物来养育既柔嫩又无经验的人，但这食物并非不经辛劳而得，无疑是为他们的益处，免得他们再度以犯罪为乐，如果加于罪的劳苦不曾唤起无罪。既然现在已不再是照料一个乐园，而是要耕种整个大地，便向人提供更强韧的肉类食物，为文化的益处，给人体增加一些活力。所有这些，如我所说，都是出于恩宠和天主的安排：以免最强韧的食物给得太少，不能维持人，使他们在劳动中衰弱；或者更柔嫩的食物不应过于丰盛，使他们被超出力量的重负压倒，无法承受。但随后颁布的法律规定了肉食的区分：它给予并准许使用某些动物，视为洁净的；禁止某些动物，视为不洁的，且使吃它们的人受污染。此外，它给洁净动物这一特征：反刍和分蹄的是洁净的；两者都不做的便是不洁的。同样，在鱼类中，法律说那些有鳞有鳍的是洁净的，而其他的是不洁的。此外，它在飞禽中设立了区别，规定了哪些应被视为可憎之物或洁净之物。因此，法律规定了\Emph{运用}极大的细致，在那些古时安排曾归于一类祝福的动物中做了区分。那么，我们该说什么？难道这些动物因此是不洁的吗？但说它们不洁净，除了法律将它们从食物用途中分离出来外，还能\Emph{说}什么呢？再者，我们刚才所说的又是什么？那么，天主是不洁之物的设立者；被造事物的过犯将回弹到创造者身上，祂没有创造它们为洁净的。说这样的话无疑是极端的、过度的愚昧：这是指控天主创造了不洁之物，并归咎于神圣的尊威制造了可憎之物的罪责，尤其当它们曾被宣布为\BibleQuote{非常好}\BibleRef{创 1:31}，且因着善从天主自己获得了祝福\Quote{要孳生繁殖}。此外，它们因造物主的命令被保留在诺厄方舟中，为了它们的后代，好使它们被保存后显明为必需的；而作为必需的，便显明为善的，尽管即便如此也附加了区别。但是，即使那时，若因自身的污染而需要消灭，那些本身不洁的受造物的创造本可以被完全废除。\par

% structure: p0006 source=h2 target=section reason=relative-heading-rank; base=h2

\section{第三章 论据——不可责难不洁的动物，以免责难落于其创造者；但当一个无理性的动物因任何理由被拒绝时，倒是那在有理性的人身上的那同一件事应受谴责；因此，在动物身上描述了人的品性、行为和意志。}

那么，那法律——我已藉宗徒的权威证明它是属神的——必须在多大程度上被以属神的方式领受，以使法律的神圣而确定的概念得以实现呢？首先，我们必须相信，凡由天主所规定的事，藉祂创造的权威本身，就是洁净且被净化的；也不可加以指责，以免这指责落在其作者身上。此外，法律\Emph{也}赐予以色列子民为了这个目的：使他们能藉以获得益处，并回归那些他们从祖先领受、却因与野蛮民族交往而在埃及腐化了的美德习尚。最后，刻在石板上的十诫并非教导什么新事，而是提醒他们已被泯灭之事——好使他们内里那已沉睡的义德，藉法律的\OriginalTerm{灵感}{afflatus}，如同被\Emph{闷住}的火一般重新燃起。但他们能从这一认识中获益：那些即使在走兽中也被法律所谴责的恶习，在人身上尤其应当避免。因为当一个无理性的动物因某种原因而被拒绝时，其实正是在理性的人身上被谴责的那同一件事。如果动物中按其本性具有的任何东西被指为不洁，那么当在人身上发现那违反其本性的同一件事时，就更应受到责备。因此，为使人得以净化，牲畜受了谴责——也就是说，那些具有同样恶习的人，也可能被看作与禽兽同等。由此得出，动物并没有因其创造者的作为而被定罪；反而人可从禽兽中得到教训，回归自己受造时那无玷的本性。因为我们必须考虑主如何区分洁净与不洁净。它说，洁净的动物既\OriginalTerm{反刍}{chew the cud}又\OriginalTerm{分蹄}{divide the hoof}；不洁净的要么两者都不做，要么只做其一。这一切都由一位工匠造成，而制造它们的那一位亲自祝福了它们。因此我认为两者的创造都是洁净的，因为创造它们的那一位是圣的，而被造之物就其受造而言并非有过。因为惯常背负罪责的并非本性，而是乖谬的意志。那么，情况如何呢？在动物身上，所描绘的是人的品性、行为和意志。它们若反刍，就是洁净的；也就是说，若他们常将天主的诫命放在口中当作食物。它们若分蹄，就是指他们以无罪的坚定脚步，行走在义德及一切生活美德的道路上。因为那些将蹄分成两瓣的动物，步履总是稳健有力；一部分蹄的滑跌倾向被另一部分的坚固所支撑，从而保持在坚实的脚步中。因此，两者都不做的，就是不洁净的；他们的行径在美德上既不坚定，也不像反刍那样消化天主诫命的食物。而那些只做其中一件的，本身也不洁净，因为他们缺少另一件，不完全具备两者。那些两者都做的，如信徒，就是洁净的；只做一件的，如犹太人和异端者，就是有瑕疵的；两者都不做的，如外邦人，就是不洁净的。如此，在动物身上，藉着法律，仿佛建立了一面人类生活的镜子，人可在其中思想惩罚的图像；这样，人身上一切违背本性而行的恶事，就可更受谴责，因为即使是那些在禽兽中自然存在的事，也在它们身上受责备。因为就鱼而言，鳞片的粗糙被视为构成其洁净；粗糙、崎岖、未打磨、实在和庄重的举止在人与上受赞赏；而那些无鳞的是不洁净的，因为轻浮、多变、不忠、柔弱的举止不受认可。此外，当法律说\BibleQuote{你不可吃骆驼？}时，是什么意思呢？——岂非藉那动物的榜样，谴责一种软弱无力且因罪行而扭曲的生活？或者当它禁止取猪为食物时呢？它无疑谴责一种污秽肮脏、以恶习的垃圾为乐的生活，将其至善不放在心灵的高尚，而只放在肉体上。或者当它禁止野兔时？它斥责那些变成女人的男人。谁会用鼬鼠的身体作食物呢？但在此它责备偷窃。谁会吃蜥蜴呢？但它憎恶漫无目的的任性生活。谁吃水蜥呢？但它痛恨心灵的污点。谁会吃鹰、鸢、鹫呢？但它憎恨靠罪行生活的掠夺者和强暴之人。谁会吃秃鹫呢？但它咒诅那些藉他人死亡寻求战利品的人。或者谁会吃乌鸦呢？但它控诉狡猾的意志。此外，当它禁止麻雀时，它谴责不节制；禁止猫头鹰时，它憎恨那些逃避真理之光的人；禁止天鹅时，它憎恨昂首骄傲的人；禁止海鸥时，它憎恨多言不节的舌头；禁止蝙蝠时，它憎恨那些寻求黑夜与错误的黑暗的人。如此，法律在动物中咒诅这些以及类似的事，这些事在动物身上确实无可指责，因为它们生来如此；但在人身上受责备，因为是人违反其本性，不是由于受造，而是由于错误而寻求的。\par

% structure: p0008 source=h2 target=section reason=relative-heading-rank; base=h2

\section{第四章．论据——此外，还有另一个理由禁止犹太人多种肉类；即为了约束百姓的不节制，并使他们事奉唯一的天主。}

于是，在这些已枚举的考虑之外，还添加了其他原因，使得许多肉类被禁止给犹太人；并且为了达到这一目的，许多东西被称为不洁，并非因其本身被定罪，而是为使犹太人受约束，只侍奉唯一的天主；因为节俭和饮食适度正适合为此目的而被拣选的人。而这种适度常被发现近乎宗教，更有甚者，可谓与之相关而相近；因为奢侈敌视圣洁。当端庄不被顾惜时，宗教又怎能被它放过呢？奢侈不存有天主的敬畏；因为当享乐催促它时，它被带向欲望的唯一胆大妄为：缰绳被松开，它无度地增加开销，仿佛这是它的食物，超出其财产及其端庄；或如从山顶急冲而下的湍流，不仅越过阻挡它的一切，而且带着那些障碍物去摧毁其他事物。因此寻求这些补救措施以抑制百姓的无节制，好使随着奢侈的减少，德行能得以增加。因为他们不配被禁止享用各种肉类的图画吗？这些人竟敢将埃及最卑贱的肉类置于玛纳的天上宴席之上，偏爱敌人和主人的多汁肉食胜过他们的自由。如果那更令人渴望和自由的食品如此不讨他们喜悦，那么他们实在配得他们所贪恋的奴役来纵容他们。\par

% structure: p0010 source=h2 target=section reason=relative-heading-rank; base=h2

\section{第五章 论证——但这些影子和预像有其限度；因为后来，当法律的终末基督来到时，宗徒便说凡洁净的人一切都是洁净的，真正的圣洁之食是正直的信德和无玷的良心。}

现在，那古时使用这些影子和预像的时期已经过去，那时要禁食那些被创造所称赞却被法律所禁止的肉类。但现在，法律的终末基督来了，揭示了法律的一切隐晦——所有那些古时以圣事的云雾所掩盖的事物。因为杰出的老师、天上的导师和完美真理的制定者已经来临，在他之下，终于正当地说：\BibleQuote{凡洁净的人，一切都是洁净的；但那些被玷污和不信的人，什么都不洁净，就连他们的理智和良心也被玷污了。} \BibleRef{铎 1:15} 而且在另一处：\BibleQuote{因为天主所造的一切都是好的，凡以感恩之心领受的，没有一样是可弃绝的；因为它是藉着天主圣言和祈祷而圣化的。} \BibleRef{弟前 4:4} 又在另一处：\Quote{圣神明明地说：在末世，有些人将背弃信德，听从欺诈的神和魔鬼的教诲，他们说话假善，良心被烧焦；他们禁止嫁娶，又\Emph{命令}戒食一些食物，这些食物是天主为了那些信者并认识天主的人，以感恩之心领受而创造的。} 而且在另一段：\BibleQuote{凡在市场上出售的，你们都可以吃，不必为良心而追问什么。} \BibleRef{格前 10:25} 从这些事可以清楚地看出，既然法律已完结，所有那些事都已回归其原有的福乐，我们不应回到那些关于肉类的特殊规条，这些规条曾因某种原因而设立，但如今福音的自由已将其取消，给予它们的解除。宗徒高呼：\BibleQuote{天主的国不在于吃喝，而在于正义、和平与喜乐。} \BibleRef{罗 14:17} 也在别处：\BibleQuote{食物是为肚腹，肚腹是为食物；但天主将使两者都朽坏。身体不是为了奸淫，而是为了主，主也是为了身体。} \BibleRef{格前 6:13} 天主不是被肚腹或食物所敬拜的，主说这些将要朽坏，并藉着自然律在排泄中被\Quote{清除}。因为凡藉食物敬拜主的，不过是以自己的肚腹为主。我再说，真正的、圣洁的、纯净的食物，乃是真正的信德、无玷的良心和无罪的灵魂。凡如此滋养的，也与基督一同滋养。这样的宴席是天主的客人；这些是滋养天使的盛宴，这些是殉道者所摆设的筵席。因此有法律的话：\BibleQuote{人不只靠饼生活，更靠从天主口中发出的每一句话而生活。} \BibleRef{申 8:3} 因此也有基督的话：\BibleQuote{我的食物就是承行派遣我者的旨意，并完成祂的工作。} \BibleRef{若 4:34} 因此，\BibleQuote{你们找我，不是因为看见了奇迹，而是因为吃饼吃饱了。不要为那可朽坏的食物劳碌，而要为那存留到永生的食物劳碌，这食物是人子要赐给你们的，因为父已印证了祂。} \BibleRef{若 6:26} 我再说，藉着正义、节制和其他德行，天主被敬拜。因为匝加利亚也告诉我们说：\Quote{你们若吃或喝，岂不是你们自己吃或喝吗？} ——借此表明食物或饮料不达到天主，而达到人；因为天主不是属血肉的，所以不会因血肉而喜悦；祂也不关心这些享乐，以致因我们的食物而欢喜。天主只喜悦我们的信德，只喜悦我们的无罪，只喜悦我们的真理，只喜悦我们的德行。这些不在我们的肚腹中，而在我们的灵魂中；这些是由神圣的敬畏和天上的畏惧而为我们所得，而非由地上的食物。宗徒适当地责备了这样的人，说：\BibleQuote{他们服从天使的迷信，因血肉的理智而自大；不持守基督为元首，全身由关节联络，藉着彼此相联的肢体交织并一同成长，在爱德之纽带中增长，以致于天主。} \BibleRef{哥 2:18} 却遵守这些东西：\BibleQuote{不可摸，不可尝，不可碰；这些事确实有虔诚的外表，但在克苦身体方面并无任何益处。} \BibleRef{哥 2:21} 然而，如果我们被自愿的奴役召回到那些我们已藉着圣洗而死亡的元素中，那么对于正义就毫无益处了。\par

% structure: p0012 source=h2 target=section reason=relative-heading-rank; base=h2

\section{第六章 论点——但既然在肉类上有了自由，并未允许奢侈，也未取消节制和斋戒：因为这些事十分相称于信徒——即他们应日夜向天主祈祷和感谢。}

但既然在肉类上有了自由，并非必然允许奢侈；福音待我们极为宽厚，却并未取消节制。我说，这并非为了满足肚腹，而是显明了肉类的形式：表明了什么是对的，不是要我们投入欲望的深渊，而是为法律提供理由。但没有什么比福音更能约束放荡；也没有人像\Person{基督}那样制定了反对贪食的严格法律，祂甚至称穷人为有福，饥饿口渴为有福，富人为可怜；那些顺从肚腹和口舌治理的人，他们情欲的原料永远不会缺乏，以致他们的奴役不能停止；他们认为能够尽情欲望是幸福的证明，只是他们因此所能得到的少于他们所欲望的。此外，祂甚至偏爱\Person{拉匝禄}的饥饿和疮疾，以及富人的狗，以榜样约束了救恩的毁灭者——肚腹和口舌。宗徒也说：\Quote{只要有衣有食，就应当知足，}\BibleRef{弟前 6:8}制定了节俭和节制的法律；他认为仅仅写下来好处不大，也把自己所写的作为榜样，并不无理由地补充说：\Quote{贪财是万恶的根；}\BibleRef{弟前 6:10}。因为它是跟随奢侈的。奢侈因恶习所耗尽的，贪婪便以罪恶来复原；犯罪循环重蹈，以便奢侈再次夺走贪婪所积攒的一切。然而，在这些事中，并不缺少那些自称基督徒之名的人，他们提供放荡的榜样和教导；他们的恶习甚至到了这样的地步，以致在斋戒时清早饮酒，认为饭后饮酒不合基督徒身份，除非酒在空腹和无阻的血管中直接下肚，紧随睡眠之后：因为他们认为若食物与酒混合，所饮的味道似乎不那么醇。因此，你能看到这类人，仍守斋却已醉，不是跑向酒馆，而是随身带着酒馆；若有人向他们问候，他们不是亲吻，而是干杯。饭后他们能做什么？因为食物已使他们沉醉？或者在日落时他们处于何种状态？因为日出时看他们已被酒弄得迟钝？但可憎之事不应作为我们的榜样。因为只有那些能使我们灵魂变得更好的事才应效法；虽然福音普遍赐予我们使用肉类的自由，但应理解为仅在节俭和节制的法律下赐予。因为这些事甚至十分相称于信徒——即那些将向天主祈祷和感谢的人，不仅白天，而且夜间也如此；因为如果头脑被酒肉麻痹，无法摆脱沉重的睡眠和压在胸口的负担，就不能如此。\par

% structure: p0014 source=h2 target=section reason=relative-heading-rank; base=h2

\section{第七章 论点——此外，必须谨慎，免得有人以为这自由可以大到甚至接近祭偶像之物。}

但在食物使用中必须极其谨慎防范，必须警告，免得有人认为自由允许到甚至可接近祭偶像之物。因为就天主的创造而言，各样受造物都是洁净的。但一旦献给魔鬼，只要献给偶像，它就受污染；一旦如此，它就不再属于天主，而属于偶像。当这受造物被取为食物时，它滋养那取食的人，不是为天主，而是为魔鬼，使他成为偶像的同席者，而非基督的同席者，正如犹太人所做的那样。明白了这些肉类的意义，思虑了法律的劝诫，认识了福音恩宠的仁慈，遵守了节制的严格，拒绝了祭偶像之物的玷污，我们这些在一切事上持守真理规则的人，应当藉着祂的子、我们的主\Person{耶稣基督}感谢天主，愿赞美、尊荣和荣耀归于祂，直到永远。阿们。一封由罗马司铎\Person{诺瓦西安}以罗马神职名义写给\Person{西彼廉}的信函翻译在此卷（第30封）第308页。\par

\SourceRecord{Translated by Robert Ernest Wallis.；Ante-Nicene Fathers；Vol. 5.；Edited by Alexander Roberts, James Donaldson, and A. Cleveland Coxe.；Buffalo, NY: Christian Literature Publishing Co.,；1886.；<http://www.newadvent.org/fathers/0512.htm>.}
